% !TEX root = ../nc-quadric-avatar.tex
% 2026-11-23 / 全基底上の有限表示とcanonicalな相対認識。
\subsection{全H2上の通常のsquare Morita代数}
\label{subsec:qa-universal-h2-square-morita81}

本項では古典型を含む全ての $\HHstack$ を扱う．
normal foundation の選択は有限 \'etale 二重被覆
$\mathfrak H_2^{\rm rul}\to\HHstack$ 上で行い，
実際の対象の代表とその Hom 束はさらに fppf 局所で取る．
そのような chart で
\[
 (E_0,E_1,E_2,E_3)=(\mathcal O,Y,X,M),
 \qquad V_i=\operatorname{Hom}(\mathcal O,E_i),
 \qquad (r_0,r_1,r_2,r_3)=(1,2,2,4)
\]
と置く．ここで $Y$ は normal 第四項の evaluation kernel であり，
$M$ は内在的な第一軌道対象である．
四対象を大域に順序付けたり，基底線束を自明化したりする必要はない．

\begin{proposition}[canonical な階数八十一の通常代数]
\label{prop:qa-universal-h2-square-morita81}
局所対象
\[
 T^{\rm can}=\bigoplus_{i=0}^3 E_i\otimes V_i^\vee
 \tag{T81}\label{eq:qa-h2-canonical-amplified-generator}
\]
は，基底線束による代表変更と二つの ruling の交換に自然に整合し，
全 $\HHstack$ の相対 heart の対象へ降下する．
その End は通常の有限局所自由代数
\[
 \mathcal A_{81}=\operatorname{End}(T^{\rm can}),
 \qquad \operatorname{rk}\mathcal A_{81}=81
\]
を定める．各相対 $H_2$ 族への引戻しで $T^{\rm can}$ は基底平坦であり，
共通 fppf 局所 chart 上の canonical compact 圏の tilting generator となる．
従って，実際の局所 Morita 同値
\[
 \operatorname{Perf}\mathcal A_{81}\simeq\mathcal C^{\rm perf}
\]
が得られる．$\mathcal O$ の直和因子に対応する大域冪等元 $e_0$ を持ち，
$e_0\mathcal A_{81}e_0=\OO$，$e_0\mathcal A_{81}\mapsto\mathcal O$ である．

代数，Morita datum，Hom 束および全積は任意の scheme 基底変更と可換する．
ただし相対 heart の主張は本稿の局所有限型 $k$-scheme の site で述べ，
大域の unbounded 導来圏と局所 compact 圏の descent の同値を新たに仮定しない．
\end{proposition}
\begin{proof}
まず局所 normal foundation を $(\mathcal O,X,M,Z)$ とする．
論文Bの相対 ruling mutation によって
\[
 Y=\operatorname{Cone}
 \bigl(\operatorname{Hom}(M,Z)\otimes M\to Z\bigr)[-1]
\]
を取り，$(\mathcal O,Y,X,M)$ は相対 full strong square 四項となる．
これは一般の $U_0$ 上でも成立する有限 Hom の計算である．
$\mathfrak F_{\HH}$ 上では $Y$ は他方の実際の normalized ruling であるので，
四対象はいずれも基底平坦な heart 対象となる．
系~\ref{cor:qa-h2-square-foundation-open} により二積は全射である．

幾何点のこの square は既存の square AS 分類の有限基礎でもある．
実際，論文Aの \cite[\texttt{thm:normal-helix}]{QuadricIntrinsic2026}
の normal helix を一項前から始めると
\[
 (H_0,H_1,H_2,H_3,H_4)=(t^{-1}Z,\mathcal O,X,M,Z)
\]
となる．同論文の \cite[\texttt{cor:normal-cubic-defect}]{QuadricIntrinsic2026}
は最小単純解消を与え，合成階数が四なら cubic，三なら最小十本矢型である．
定理~\ref{thm:qa-as10-inverse-all-targets} と
補題~\ref{lem:qa-as10-all-index-mutual-inverse} の逆操作は
\[
 \operatorname{coker}
 (t^{-1}Z\to\operatorname{Hom}(t^{-1}Z,\mathcal O)^\vee\otimes\mathcal O)
   \simeq Y
\]
を与え，この square を回収する．古典型では通常の
$(\OO,\OO(1,0),\OO(0,1),\OO(1,1))$ である．
従って局所有限代数 $\Lambda=\operatorname{End}(\bigoplus_iE_i)$ は，
補題~\ref{lem:qa-classified-fiber-finite-square-recognition} の範囲に入る．
ここで大域的な曲線 avatar が存在するとは仮定していない．

\emph{頂点線束の消去．}
二つの局所代表の比較で $E_i'=E_i\otimes L_i$ なら，
pointed な $\mathcal O$ の比較の下で $V_i'=V_i\otimes L_i$ である．
従って自然な縮約により
\[
 E_i'\otimes(V_i')^\vee
 \simeq E_i\otimes L_i\otimes V_i^\vee\otimes L_i^\vee
 \simeq E_i\otimes V_i^\vee
\]
を得る．比較射を単元倍した際，対象上の単元と Hom 双対上の逆単元は
厳密に相殺される．縮約は tensor の結合と整合するので，三重重なりの
scalar cocycle も消える．これは有限 square 表示の頂点トーラス
$K\simeq\Gm^3$ の全作用を消すことに等しい．

pointed exact 同値は二つの ruling の非順序対と $M$ を保つ．
ruling の交換は上の直和の二つの中間項だけを入れ替える．
有限直和の自然な置換は二回で恒等となり，既存の ruling 交換の coherent
$S_2$ 作用と整合する．よって $T^{\rm can}$ は ruling を選ばない
$\HHstack$ 上へ降下する．$\mathcal O$ 項と $M\otimes V_3^\vee$ 項は
個別に降下し，中間二項についてはその\emph{和}が降下する．

同じ議論は代数を直接構成しても確認できる．局所の $(j,i)$ corner は
\[
 \operatorname{Hom}(E_i,E_j)\otimes V_j^\vee\otimes V_i
\]
であり，頂点スカラーの指標
$(\lambda_j/\lambda_i)\lambda_j^{-1}\lambda_i$ は一である．
合成は中間の $V_j\otimes V_j^\vee$ の縮約なので，この降下は全積を保つ．

\emph{階数とtilting．}
局所 square の Hom 階数行列を
\[
 G=\begin{pmatrix}1&2&2&4\\0&1&0&2\\0&0&1&2\\0&0&0&1\end{pmatrix},
 \qquad r=(1,2,2,4)^{\mathsf T}
\]
とすれば，End の階数は
\[
 r^{\mathsf T}Gr
 =(1+4+4+16)+(4+4+16+16+16)=81
\]
である．各 $V_i$ を局所自明化すれば $T^{\rm can}$ は $E_i$ を
$r_i$ 回ずつ増幅した四項の直和であり，full strong 性と tilting が従う．
$r_0=1$ と単位 $V_0=\operatorname{End}(\mathcal O)=\OO$ によって
指定冪等元の同定も得る．
局所的な compact Morita datum の比較においても，前述の scalar cancellation
と $S_2$ の置換が cocycle を与える．
全 Hom 束は perfect-fiber 判定による有限局所自由束であり，
この対角・tensor・縮約の構成は任意基底変更に保たれる．
\end{proof}

\begin{remark}[大域の対角代数]
\label{rem:qa-h2-rank81-diagonal}
この構成は，二つの ruling を大域にラベル付けした四つの分裂頂点を
要求しない．ruling 選択被覆を $\pi:\mathfrak H_2^{\rm rul}\to\HHstack$ とすると，
相対的な対角代数は
\[
 \mathcal D=\OO\oplus\pi_*\mathcal D_{\rm rul}\oplus\mathcal D_M
\]
と書ける．$\mathcal D_{\rm rul}$ は選択被覆上で局所的に
$\operatorname{End}(V_X^\vee)$ となる次数二の Azumaya 代数，
$\mathcal D_M$ は局所的に $\operatorname{End}(V_3^\vee)$ となる
次数四の Azumaya 代数である．それらを大域的な行列代数と仮定しない．
対角の階数は $1+8+16=25$ である．
厳密に前向きの corner の和を $\mathcal J$ とすると，
\[
 \mathcal A_{81}=\mathcal D\oplus\mathcal J,
 \qquad \mathcal J^{\otimes_{\mathcal D}3}=0.
\]
従って大域の separable 対角を持つ有限 bar 解消を適用できる．
\end{remark}

\begin{theorem}[全H2のcanonical完成とpointed tails]
\label{thm:qa-universal-h2-rank81-completion}
上の通常代数に対して
\[
 \mathcal C_{81}
 =\bigoplus_{n\ge0}
 H^0\left(\Theta_{\mathcal A_{81}}^{
                   \otimes_{\mathcal A_{81}}^{\mathbf L}n}\right)
\]
を取る．全 tensor 冪は cohomological degree 零に集中し，
\[
 \operatorname{rk}(\mathcal C_{81})_n=81(2n+1)^2
 \qquad(n\ge0).
\]
全次数と全積は任意の scheme 基底変更と可換する．
零頂点 corner は canonical な反 Serre 環であり，
\[
 e_0\mathcal C_{81}e_0
   =\bigoplus_{n\ge0}\operatorname{Hom}(\mathcal O,t^n\mathcal O),
 \qquad
 \operatorname{rk}e_0(\mathcal C_{81})_ne_0=(2n+1)^2.
\]

本稿の site の任意の affine test $S\to\HHstack$ 上で，
$(\mathcal C_{81})_S$ は左右 Noetherian であり，実際の pointed exact 同値
\[
 \bigl(\qgr(\mathcal C_{81})_S,\,
       \pi e_0(\mathcal C_{81})_S,\,(1)\bigr)
   \simeq (\mathcal C_S,\mathcal O_S,t)
 \tag{QT81}\label{eq:qa-h2-rank81-pointed-tails}
\]
が存在する．この比較は通常の module-object 基底変更と coherent に整合する．
\end{theorem}
\begin{proof}
前注意の finite separable diagonal 解消により
命題~\ref{prop:qa-relative-serre-tensor-base-change} が適用できる．
局所 square の Morita 表示では，その第 $n$ tensor 冪は
\[
 \mathbf R\operatorname{Hom}(T^{\rm can},t^nT^{\rm can})
\]
に対応する．幾何 fiber の全整数 square helix の前向き Ext 消滅によって，
$n=0$ を含む全 $n\ge0$ で次数零に集中する．
perfect-fiber 判定は全基底上の局所自由性と任意基底変更を与える．
積は相対 Serre transport と合成によるものなので，同じ比較を保つ．

数値 Serre 行列を $S_G=G^{-1}G^{\mathsf T}$ とすると，fiber の階数は
$r^{\mathsf T}GS_G^{-n}r$ である．$(S_G-I)^3=0$ を用いた直接計算は
\[
 r^{\mathsf T}GS_G^{-n}r=81(2n+1)^2
\]
を与える．これは前向き Ext 消滅を別に証明した後の階数計算である．
$e_0$ が $\mathcal O$ を保つため，corner は全積を含めて
表示した canonical 反 Serre 環に一致する．

\emph{左右Noether性．}
まず各 $V_i$ と四頂点が分裂する局所 chart を取る．
$B=C(\Lambda)$ と書くと，$(\mathcal C_{81})_n$ は
$e_jB_ne_i$ を $r_j\times r_i$ 個に増幅した block 行列である．
各頂点の $r_i$ コピーから一つずつを選ぶ零次数冪等元を $f$ とすれば
\[
 f\mathcal C_{81}f\simeq B.
\]
零次数の同一頂点内の matrix units によって
$\mathcal C_{81}f\mathcal C_{81}=\mathcal C_{81}$ である．
従ってこれは通常の degree-zero graded Morita 同値である．
局所 $\Lambda$ の全幾何 fiber は分類済み square 型であり，
補題~\ref{lem:qa-classified-fiber-finite-square-recognition}，
\ref{lem:qa-square-foundation-opposite}，
\ref{lem:qa-serre-completion-opposite} によって $B$ は左右 Noetherian である．
この性質は通常の Morita 同値で保たれ，中央の faithfully flat 基底変更で
降下する．よって全ての指定された affine test 上で左右 Noether 性を得る．
ここで用いた full idempotent $f$ は各増幅頂点からの一コピーの和であり，
$\mathcal O$ だけの冪等元 $e_0$ とは区別する．

\emph{実際のheart．}
$T^{\rm can}$ は基底平坦 compact な対象の有限直和で，
$\mathcal O$ を直和因子に含む．各項から十分前向きの
$t^m\mathcal O$ への高次 Ext は既に消えている．
従って命題~\ref{prop:qa-relative-flat-foundation-az} により
$(T^{\rm can},t)$ は相対 AZ ample である．
実際の軌道環は，全積を含めて
\[
 \bigoplus_{n\ge0}\operatorname{Hom}(T^{\rm can},t^nT^{\rm can})
    =\mathcal C_{81}
\]
である．一般の有限 End 環に対する AZ 再構成により
式~\eqref{eq:qa-h2-rank81-pointed-tails} を得る．
これは $e_0$ の corner だけから graded Morita 同値を導いたものではない．
指定対象は $T^{\rm can}$ の実際の直和因子として識別される．
基底変更との比較は，標準射影のシフトを $t$ の軌道対象へ送り，
有限表示の余核へ延長する構成である．したがって低次数 section の
非平坦基底変更を仮定せずに，通常の module-object 引戻しとの
coherent な比較を得る．
\end{proof}

\begin{remark}[この通常代数の役割]
この結果は，古典型を含む全ての内在的 $H_2$ 族に，
指定構造対象を保持する通常の有限 Morita 代数とその canonical 全 Serre 完成を
付随させる．四つの rank-one 対象の大域的な代表や，ruling の順序は不要である．
一方，任意の AS/QP presentation stack を同定するには，
その presentation の相対平坦性，実際の全整数積，頂点線束と Serre transport，
および全同型層を比較する別の証明が必要である．
また任意の Noetherian $k$-基底へ相対 heart の定義を拡張したものとの
同一視や，その意味での strong Noether 性は本定理に含めない．
\end{remark}
